%!TEX root = main.tex

\section{Evaluation}
\label{sec:evaluation}


\subsection{Experimental Setup}
\label{sub:setup}


For our performance comparison, we use OMNET++4.6/INET3.3, a simulation framework specifically aimed at networking algorithms.
%
This tool provides models of many protocols of the network stack such as MAC, 802.11g wireless communication and energy consumption ~\cite{10.4108/ICST.SIMUTOOLS2009.5585}.
%
In addition, it contains models to describe and deal with the mobility aspect.
%
In the following, we describe the simulation environment and the parameters used. 

\subsubsection*{Communication Scenario}

Our experiments are inspired by the scenario discussed in Section~\ref{sec:scenarios}.
%
In this scenario, attendees have a mobile device with IEEE 802.11 capability.
%
Many participants might attempt to broadcast messages at nearly the same time using a MANET.
%
Taking this into consideration, in our performance comparison  we set the number of broadcast messages to 300 and the interval among messages to the values $0.01$ and $0.5$ seconds, which are less than the speed of moving nodes.
%
These two values allow us to evaluate the protocols under heavy and light network load.
%
Since some algorithms build a topology to drive the flooding operation, there is a warm-up period before the first message is broadcasted that protocols can use to create any initial structure.
%
Finally, in our evaluation we stop each experiment a fixed amount of time after the beginning of the last flooding operation.

\subsubsection*{Density}

\subsubsection*{Mobility}

In our evaluation, we use the \textit{truncated-levy walk} (LW) mobility model, as it characterizes the mobility traces of humans carrying wireless devices~\cite{Rhee:2011}.
%
This model generates traces composed of frequent short walks and occasional long rides to distant locations.
%
In our context, we define two scenarios: with mobility at walking speed and without mobility.
%
We set the parameter of our LW model in such a way that the average walking speed is close to $1.4m/s$, a value reported in the literature~\cite{Browning390} as an estimation of the preferred walking speed.
%
To do so, we use the following values:

\begin{table}[ht]
	\footnotesize
	\centering
	
	\begin{tabular}{lcp{3.5cm}}
		\textbf{Parameter} & \textbf{Value} & \textbf{Description} \\
		
		\midrule
		FL\_EXP & -3.9 & Exponent of the flight length distribution \\
		FL\_MAX & 50 & Maximum value of the flight length distribution \\
		WT\_EXP & -1.8 & Exponent of the waiting time distribution \\
		WT\_MAX & 100 & Maximum value of the waiting time distribution \\
	\end{tabular}
	
	\caption{Configuration of the LW mobility model.}
	\label{tbl:mobility-parameters}
\end{table}

Afterwards, we check the traces to guarantee that the average speeds are indeed close to $1.4m/s$, and the topology remains connected all the time.

\subsubsection*{Energy Consumption}


\subsubsection*{Algorithm specific parameters}

The behavior of each protocol depends on the values of its parameters.
%
To configure these parameters in a meaningful way, we choose values previously reported in the literature.
%
For instance, we use 0.5 seconds as the \textit{hello time} for both \textit{MPR} and \textit{node-degree-based} CDS.
%
In the case of \textit{ABBA}, we use 0.5 seconds as the maximum waiting time.
%
Finally, for \textit{ProbFlood} we assign the values $k=15$, $\sigma=1$, and $\textit{hello time} = 0.5$.

\subsection{Performance Metrics}

\paragraph{\textbf{Saved rebroadcasts}}
number of messages that the nodes do not send thanks to the use of a particular broadcast algorithm.
%
More formally, let $k$ be the number of observed retransmissions during the flooding of a message using algorithm $a$; and $n$ the number of nodes in the network; then, we save $n - k$ rebroadcast by using $a$ instead of \textit{simple flooding}.
%
This metric is useful to measure the efficiency in terms of both network and energy usage. 

\paragraph{\textbf{Duplicated messages}}
 number of extra times a message is received.
%
That is, in a perfect situation a node only need to receive a message once.
%
But with most algorithms the same messages is usually received multiple times.

\paragraph{\textbf{Coverage}}
 percentage of nodes in the network that receive a message.


\paragraph{\textbf{Broadcast time}} time required for a flood operation to reach full coverage.
%
It is only measured when all the participants receive the message because this is the only way of comparing protocols with different coverage. 

\paragraph{\textbf{Energy consumption}}
 energy consumed for each node during the experiment.




\begin{figure*}[!ht]
	\centering
	\includegraphics[page=1, width=0.95\textwidth]{fig/results-static.pdf}
	\vspace{-0.5cm}
	\caption{Power consumption when there is no mobility}
	\label{fig:power-consumption-static}
\end{figure*}


\begin{figure*}[!ht]
	\centering
	\includegraphics[page=2, width=0.95\textwidth]{fig/results-static.pdf}
	\vspace{-0.7cm}
	\includegraphics[page=4, width=0.95\textwidth]{fig/results-static.pdf}
	\vspace{-0.0cm}
	\caption{Network usage when there is no mobility: a) number of saved rebroadcasts, b) number of duplicated messages}
	\label{fig:network-usage-static}
\end{figure*}


\begin{figure*}[!ht]
	\centering
	\includegraphics[page=1, width=0.95\textwidth]{fig/results-mobility.pdf}
	\vspace{-0.5cm}
	\caption{Power consumption with mobility}
	\label{fig:power-consumption-mobility}
\end{figure*}

\begin{figure*}[!ht]
	\centering
	\includegraphics[page=2, width=0.95\textwidth]{fig/results-mobility.pdf}
	\vspace{-0.7cm}
	\includegraphics[page=4, width=0.95\textwidth]{fig/results-mobility.pdf}
	\vspace{-0.0cm}
	\caption{Network usage with mobility: a) number of saved rebroadcasts, b) number of duplicated messages}
	\label{fig:network-usage-mobility}
\end{figure*}